\documentclass[10pt,a4paper]{amsart}
\usepackage[margin=19mm]{geometry}
\usepackage{amsmath,amssymb,mathtools}
\usepackage[scr=rsfs,cal=euler]{mathalfa}
\usepackage{xfrac}
\usepackage[unicode,hidelinks,bookmarks=false]{hyperref}
\newcommand{\ie}{i.e.\ }
\newcommand{\cf}{cf.\ }
\newcommand{\resp}{respectively\ }
\newcommand{\Subsection}{\subsection}
\providecommand{\nobreakdash}{\nobreak\hbox{-}}
\setlength{\emergencystretch}{2em}
\multlinegap0pt
\newcommand{\pG}{\partial{\Gamma}}
\newcommand{\dirac}[1]{{\mathrm{Dirac}}{(#1)}}



\newcommand{\N}{\mathbb{N}}
\newcommand{\Z}{\mathbb{Z}}
\newcommand{\bbP}{\mathbb{P}}
\newcommand{\R}{\mathbb{R}}
\newcommand\bbN{\mathbb{N}}
\newcommand{\rme}{\mathrm{e}}
\newcommand{\rmd}{\mathrm{d}}

\newcommand{\fkd}{\mathfrak{d}}
\newcommand{\fkB}{\mathfrak{B}}
\newcommand{\bfE}{\mathbf{E}}
\newcommand{\bfP}{\mathbf{P}}

\newcommand{\clM}{\mathcal{M}}
\newcommand{\clP}{\mathcal{P}}
\newcommand{\clV}{\mathcal{V}}
\newcommand{\clZ}{\mathcal{Z}}
\newcommand{\clW}{\mathcal{W}}


\DeclareMathOperator{\BRW}{BRW}
\DeclareMathOperator{\supp}{supp}
\DeclareMathOperator{\con}{con}
\DeclareMathOperator{\Cay}{Cay}
\DeclareMathOperator{\sd}{sd}


%%%%%%%%%%%%%%%%%%%%%%%%%%%%%%%%%%%%%%%%%%%%%%%%%%%%%%%%%%%%%%%%%%%%%%%%%%%%%%%%%%%%
%Better widebar

\makeatletter
\let\save@mathaccent\mathaccent
\newcommand*\if@single[3]{%
  \setbox0\hbox{${\mathaccent"0362{#1}}^H$}%
  \setbox2\hbox{${\mathaccent"0362{\kern0pt#1}}^H$}%
  \ifdim\ht0=\ht2 #3\else #2\fi
  }
%The bar will be moved to the right by a half of \macc@kerna, which is computed by amsmath:
\newcommand*\rel@kern[1]{\kern#1\dimexpr\macc@kerna}
%If there's a superscript following the bar, then no negative kern may follow the bar;
%an additional {} makes sure that the superscript is high enough in this case:
\newcommand*\widebar[1]{\@ifnextchar^{{\wide@bar{#1}{0}}}{\wide@bar{#1}{1}}}
%Use a separate algorithm for single symbols:
\newcommand*\wide@bar[2]{\if@single{#1}{\wide@bar@{#1}{#2}{1}}{\wide@bar@{#1}{#2}{2}}}
\newcommand*\wide@bar@[3]{%
  \begingroup
  \def\mathaccent##1##2{%
%Enable nesting of accents:
    \let\mathaccent\save@mathaccent
%If there's more than a single symbol, use the first character instead (see below):
    \if#32 \let\macc@nucleus\first@char \fi
%Determine the italic correction:
    \setbox\z@\hbox{$\macc@style{\macc@nucleus}_{}$}%
    \setbox\tw@\hbox{$\macc@style{\macc@nucleus}{}_{}$}%
    \dimen@\wd\tw@
    \advance\dimen@-\wd\z@
%Now \dimen@ is the italic correction of the symbol.
    \divide\dimen@ 3
    \@tempdima\wd\tw@
    \advance\@tempdima-\scriptspace
%Now \@tempdima is the width of the symbol.
    \divide\@tempdima 10
    \advance\dimen@-\@tempdima
%Now \dimen@ = (italic correction / 3) - (Breite / 10)
    \ifdim\dimen@>\z@ \dimen@0pt\fi
%The bar will be shortened in the case \dimen@<0 !
    \rel@kern{0.6}\kern-\dimen@
    \if#31
      \overline{\rel@kern{-0.6}\kern\dimen@\macc@nucleus\rel@kern{0.4}\kern\dimen@}%
      \advance\dimen@0.4\dimexpr\macc@kerna
%Place the combined final kern (-\dimen@) if it is >0 or if a superscript follows:
      \let\final@kern#2%
      \ifdim\dimen@<\z@ \let\final@kern1\fi
      \if\final@kern1 \kern-\dimen@\fi
    \else
      \overline{\rel@kern{-0.6}\kern\dimen@#1}%
    \fi
  }%
  \macc@depth\@ne
  \let\math@bgroup\@empty \let\math@egroup\macc@set@skewchar
  \mathsurround\z@ \frozen@everymath{\mathgroup\macc@group\relax}%
  \macc@set@skewchar\relax
  \let\mathaccentV\macc@nested@a
%The following initialises \macc@kerna and calls \mathaccent:
  \if#31
    \macc@nested@a\relax111{#1}%
  \else
%If the argument consists of more than one symbol, and if the first token is
%a letter, use that letter for the computations:
    \def\gobble@till@marker##1\endmarker{}%
    \futurelet\first@char\gobble@till@marker#1\endmarker
    \ifcat\noexpand\first@char A\else
      \def\first@char{}%
    \fi
    \macc@nested@a\relax111{\first@char}%
  \fi
  \endgroup
}
\makeatother

\let\oldbar\bar
\renewcommand*{\bar}[1]{{\mathchoice{\widebar{#1}}{\widebar{#1}}{\widebar{#1}}{\oldbar{#1}}}}
%%%%%%%%%%%%%%%
%~ This fixes spacing issues around \left( and \right) and other similar symbols
\let\originalleft\left 
\let\originalright\right
\renewcommand{\left}{\mathopen{}\mathclose\bgroup\originalleft}
\renewcommand{\right}{\aftergroup\egroup\originalright}
%%%%%%%%
\let\oldin\in
\renewcommand{\in}{\mathchoice{\oldin}{\oldin}{\,\oldin\,}{\,\oldin\,}}

\let\oldge\ge
\renewcommand{\ge}{\mathchoice{\oldge}{\oldge}{\,\oldge\,}{\,\oldge\,}}

\let\oldle\le
\renewcommand{\le}{\mathchoice{\oldle}{\oldle}{\,\oldle\,}{\,\oldle\,}}

\renewcommand*{\to}{\mathchoice{\longrightarrow}{\rightarrow}{\rightarrow}{\rightarrow}}

\let\oldmapsto\mapsto
\renewcommand*{\mapsto}{\mathchoice{\longmapsto}{\oldmapsto}{\oldmapsto}{\oldmapsto}}
%%%%%%%%%%%%%%%%%%%%%%%%%%%%%%%%%%%%%%%%%%%%%%%%%%%%%%%%%%%%%%%%%%%%%%%%%%%%%%%%%%%%


\newcommand*{\mk}{\mkern -1mu}
\newcommand*{\Mk}{\mkern -2mu}
\newcommand*{\mK}{\mkern 1mu}
\newcommand*{\MK}{\mkern 2mu}

\hypersetup{urlcolor=purple, linkcolor=blue, citecolor=red}

\newcommand*{\relabel}{\renewcommand{\labelenumi}{(\theenumi)}}
\newcommand*{\romanenumi}{\renewcommand*{\theenumi}{\roman{enumi}}\relabel}
\newcommand*{\Romanenumi}{\renewcommand*{\theenumi}{\Roman{enumi}}\relabel}
\newcommand*{\alphenumi}{\renewcommand*{\theenumi}{\alph{enumi}}\relabel}
\newcommand*{\Alphenumi}{\renewcommand*{\theenumi}{\Alph{enumi}}\relabel}
\let\oldhat\hat
\renewcommand*{\hat}[1]{\mathchoice{\widehat{#1}}{\widehat{#1}}{\oldhat{#1}}{\oldhat{#1}}}
\let\oldexists\exists
\renewcommand*{\exists}{\mathrel{\oldexists}}
\let\oldforall\forall
\renewcommand*{\forall}{\mathrel{\oldforall}}

\newcommand{\llbangle}{\mathopen{\langle\mkern -3.7mu\langle}}
\newcommand{\rrbangle}{\mathclose{\rangle\mkern -3.7mu\rangle}}

\newtheorem{theo}{Theorem}[section]
\newtheorem{lemm}[theo]{Lemma}
\newtheorem{prop}[theo]{Proposition}
\newtheorem{coro}[theo]{Corollary}
\theoremstyle{definition}\newtheorem{defi}[theo]{Definition}
\newtheorem{rema}[theo]{Remark}
\title[Twisted partial-shadow bounds]{The growth of the Green function for random walks and Poincaré series: original Lemma4.14 (Shadow lemma)}
\author{Matthieu Dussaule, Longmin Wang and Wenyuan Yang}
\date{Source: Annales Henri Lebesgue8(2025),1061--1107; DOI10.5802/ahl.256}
\begin{document}\maketitle
\noindent\textit{Editorial scope.} The counted unit is original Lemma4.14 (Shadow lemma) with its complete proof and the necessary source-local core: definitions of relatively hyperbolic groups, conical and parabolic limit points and transition points, weak relative Ancona inequalities, the Busemann cocycle and Green distance, quasi-conformal densities, the Patterson--Sullivan construction and the shadow definitions. The counted proof uses no result outside this excerpt except the named DWY transition-point lemma cited in the first line. Later atom/divergence results, all of Section5 and AppendixA are omitted. Literal source slips are kept unchanged; no mathematics was written or repaired.
\subsection*{Notations}
In all the paper, given two functions $f$ and $g$, we write $f\simeq g$ if the difference between $f$ and $g$ is bounded from above and below, that is $|f-g|\leq C$ for some constant $C$. We write $f\simeq_L g$ if the constant $C$ depends on a parameter $L$. Assuming further that $f$ and $g$ are positive, we write $f\asymp g$ if the ratio of $f$ and $g$ is bounded\pagebreak{} from above and below, that is $\frac{1}{C}f\leq g\leq Cf$ for some positive constant $C$. Similarly, if $C$ depends on $L$, we write $f\asymp_L g$. Also, if $f\leq Cg$, we write $f\prec g$ and $f\prec_L g$ if $C$ depends on $L$. If the dependency is not clear from the context, we will avoid these notations.

\setcounter{section}{3}
\section{Twisted Patterson--Sullivan measures}
\label{Sectionpattersonsullivan}

Since their introduction by Gromov, relatively hyperbolic groups were studied by many authors through several equivalent definitions. We will mainly use the viewpoint of Bowditch~\cite{Bowditch} and Gerasimov--Potyagailo~\cite{GePoJEMS,GePoGGD,GP16} in the sequel. Consider a finitely generated group $\Gamma$ acting properly via isometries on a proper geodesic Gromov hyperbolic space~$X$. Define the \emph{limit set} $\Lambda_\Gamma$ as the closure of~$\Gamma$ in the Gromov boundary $\partial X$ of $X$, that is, fixing a base point $x_0$ in $X$, $\Lambda_\Gamma$ is the set of all possible limits of sequences $g_n\cdot x_0$ in $\partial X$, $g_n\in \Gamma$. The proper action of $\Gamma$ on~$X$ by isometries extends to a convergence group action on $\Lambda_\Gamma$ by homeomorphisms, which means that the induced action on the space of distinct triples is properly continuous (see~\cite[Proposition~1.11]{Bo99} for example). If $\Lambda_\Gamma$ contains at least three points, then $\Gamma$ acts minimally on $\Lambda_\Gamma$.

A loxodromic element $x\in \Gamma$ is an infinite order element with exactly two fixed points $x_-\ne x_+$ in $\Lambda_\Gamma$. Moreover, $x$ acts via North-South dynamics on $\Lambda_\Gamma$ in the sense that for any $\xi\neq x_{\pm}$, $x^{\mp n}\xi$ converges to $x_{\mp}$ as $n$ goes to $+\infty$. Then $x_+$ is called the attracting fixed point of $x$ and $x_-$ is called its repelling fixed point. The fixed points of any two loxodromic elements are either the same or disjoint. So $\Gamma$ contains infinitely many loxodromic elements with pairwise disjoint fixed points.


A point $\xi\in \Lambda_\Gamma$ is called \emph{conical} if there is a sequence $g_{n}$ of $\Gamma$ and distinct points $\xi_1,\xi_2$ in $\Lambda_\Gamma$ such that $g_{n}\xi$ converges to $\xi_1$ and $g_{n}\zeta$ converges to $\xi_2$ for all $\zeta\neq \xi$ in~$\Lambda_\Gamma$. A point $\xi\in \Lambda_\Gamma$ is called \emph{parabolic} if its stabilizer in $\Gamma$ is infinite, fixes exactly $\xi$ in $\Lambda_\Gamma$ and contains no loxodromic element. If, in addition, its stabilizer in $\Gamma$ acts co-compactly on $\Lambda_\Gamma \setminus \{\xi\}$, then $\xi$ is called \emph{bounded parabolic}. Say that the action of~$\Gamma$ on $X$ is \emph{geometrically finite} if the induced convergence group action on the limit set $\Lambda_
\Gamma$ is geometrically finite: $\Lambda_
\Gamma$ only consists of conical limit points and bounded parabolic limit points. See~\cite{Bo99} for more general facts on convergence groups.


A group $\Gamma$ is called \emph{relatively hyperbolic} with respect to a collection of subgroups~$\bbP$ if it acts geometrically finitely on a proper geodesic hyperbolic space $X$ such that the stabilizers of parabolic limit points are exactly the conjugates of the elements of~$\bbP$. Elements of $\bbP$ are called maximal parabolic subgroups. We will write $\bbP_0$ for the choice of a set of representatives of conjugacy classes of elements of $\bbP$. By~\cite[Proposition~6.15]{Bowditch}, such a set $\bbP_0$ is finite.

The limit set $\Lambda_\Gamma$ in the Gromov boundary of $X$ is called the \emph{Bowditch boundary} of $\Gamma$. By~\cite[Theorem~9.4]{Bowditch}, it is unique up to equivariant homeomorphism and in particular does not depend on the choice of a proper geodesic Gromov hyperbolic space $X$ on which $\Gamma$ acts geometrically finitely. We will write $\partial\Gamma$ for the Bowditch boundary of $\Gamma$ in the sequel. A relatively hyperbolic group is called \emph{non-elementary} if its Bowditch boundary is infinite; equivalently, $\sharp \partial\Gamma>2$.


We now fix a finite set $\bbP_0$ of representatives of conjugacy classes of maximal parabolic subgroups. When $P\in \bbP_0$ and $g\in \Gamma$, we call $gP$ a maximal parabolic coset.
\begin{defi}
Let $gP$ be a maximal parabolic coset and $\eta, L>0$ be fixed constants. A point $p$ on a geodesic $\alpha$ is called \emph{$(\eta, L)$-deep} in $gP$ if
\[
B(p, 2L)\cap \alpha\subseteq N_\eta(gP).
\]
It is called an \emph{$(\eta, L)$-transition point} if it is not \emph{$(\eta, L)$-deep} in any maximal parabolic coset $gP$.
\end{defi}

The following set of inequalities is called weak relative Ancona inequalities and will be helpful below. They extend similar inequalities proved for hyperbolic groups by Ancona~\cite{Ancona} for $r=1$, by Gou\"ezel--Lalley~\cite{GouezelLalley} on co-compact Fuchsian groups for $r\leq R$ and by Gou\"ezel~\cite{Gouezel} in full generality for $r\leq R$. The version for relatively hyperbolic groups that we use here were first proved for $r=1$ by Gekhtman--Gerasimov--Potyagailo--Yang~\cite{GGPY} and then by Dussaule--Gekhtman~\cite{DG21} for $r\leq R$.

\begin{lemm}[{\cite[Theorem~1.1]{GGPY}, \cite[Theorem~1.6]{DG21}}]
\label{weakAncona}
Let $\Gamma$ be a relatively hyperbolic group and let $\mu$ be a finitely supported admissible and symmetric probability measure on $\Gamma$. Then, for every $c,\eta,L\geq 0$ there exists $C>0$ such that for every $r\leq R$, the following holds. For every $x,y,z\in \Gamma$, if $y$ has a distance at most $c$ to an $(\eta,L)$-transition point on a geodesic from $x$ to $z$, then
\[
\frac{1}{C}G(x,y|r)G(y,z|r)\leq G(x,z|r)\leq C G(x,y|r)G(y,z|r).
\]
\end{lemm}

Note that the constant $C$ is independent of $r\in [1,R]$. In other words, the Green function is roughly multiplicative along transition points on a geodesic. In both~\cite{DG21,GGPY}, these inequalities are formulated in terms of the Floyd distance, which is an appropriate rescaling of the word distance. However, the statement for transition points directly follows from~\cite[Corollary~5.10]{GePoGGD} which relates transition points with the Floyd distance. We also refer to~\cite[Section~9]{GGPY} for more details.

We will also use the following at some point.

\begin{lemm}[{\cite[Lemma~2.14]{YangPS}}]\label{ConCharLem}
There exist universal constants $\eta, L$ with the following property. Let $\gamma$ be a geodesic ray ending at a conical point $\xi\in\pG$. Then $\gamma$ contains a unbounded sequence of $(\eta,L)$-transition points $x_n$.
\end{lemm}


In the remainder of this section, we consider a finitely generated relatively hyperbolic group $\Gamma$. When speaking of a transition point, we mean an $(\eta,L)$-transition point with $(\eta,L)$ satisfying Lemma~\ref{ConCharLem}. We also fix a finitely supported symmetric and admissible probability measure $\mu$ on $\Gamma$.

\subsection{Busemann cocycles}
Given $x,y,z \in \Gamma$, let $B_z(x, y) := d(x, z) -d(y,z)$ and $K_z(x,y|r)=\frac{G(x,z|r)}{G(y,z|r)}$. The function $B_z$ is called the Busemann function associated with the distance $d$ at $z$.

Following~\cite{BlachereBrofferio}, we define the $r$-Green distance by
\[
d_r(x,y) = - \log F_r(x,y)=-\log \frac{ G(x,y|r)}{G(e,e|r)}.
\]
Then $K_z(x,y|r)=\rme^{-[d_r(x,z)-d_r(y,z)]}$ is the exponential of the Busemann function for the $r$-Green distance. We also write $|x^{-1}y|_r=d_r(x,y)$, and $|x^{-1}y|=d(x,y)$.

Consider the distance $\fkd_r$ defined for $x,y\in \Gamma$ by
\[
\fkd_r(x,y):=\omega_\Gamma(r)\left|x^{-1}y\right| +\left|x^{-1}y\right|_r.
\]


\setcounter{theo}{5}
Define the corresponding Busemann cocycle
\begin{equation}\label{defBusemann}
\fkB_\xi(x, y;r) = \omega_\Gamma(r)B_\xi(x, y)-\log K_\xi(x,y|r)
\end{equation}



\begin{lemm}\label{BusemanEstLem}
There exists a constant $C>0$ with the following property. Let $\xi \in \pG$ be a conical point, and $x, y \in \Gamma$. There exists a neighborhood $V=V(x,y)$ of $\xi$ in $\Gamma\cup\pG$ such that for any $ z \in V \cap \Gamma$:
\begin{align*}
\left\lvert B_\xi(x, y)-B_z(x, y)\right| &\le C, \\
\left\lvert \log K_{\xi}(x, y|r)-\log K_{z}(x, y|r)\right| &\le C.
\end{align*}
\end{lemm}

\begin{proof}
The statement for $B_\xi$ is proved in~\cite[Lemma~2.20]{YangPS}. Also, by~\cite[Proposition~4.1]{DG21}, the Martin kernel $K_z(\cdot,e|r)=G(\cdot, z|r)/G(e, z|r)$ extends continuously to $K_\xi(\cdot,e|r)$ as $z$ converges to a conical limit point $\xi$. This follows from weak relative Ancona inequalities. In particular, $K_z(x,y|r)$ converges to $K_\xi(x,y|r)$ as $z$ converges to $\xi$, so the statement for $K_\xi$ also holds.
\end{proof}

As a consequence, the Busemann cocycle $\fkB_z(x,y)$ extends to a coarse cocycle $\fkB_\xi(x,y)$ at a conical point $\xi$. That is,
\[
\fkB_\xi(x,y) := \limsup_{z\to \xi} \fkB_\xi(x,y)
\]
does not depend on the choice of $z\to \xi$, up to a bounded additive error $C$ independent of $\xi$.

\subsection{Quasi-conformal densities}

A Borel measure $\mu$ on a topological Hausdorff space $T$ is \emph{regular} if $\mu(A) = \inf\{\mu(U): A \subset U, U$ is open$\}$ for any Borel set $A$ in $T$. It is called \emph{tight} if $\mu(A) =
\sup\{\mu(K): K \subset A, K$ is compact$\}$ for any Borel set $A$ in $T$. A finite Borel measure is called \emph{Radon} if it is tight and regular. It is well-known that all finite Borel measures on compact metric spaces are Radon, see~\cite[Theorem~1.1, Theorem~1.3]{Bill99}.

Denote by $\clM(\pG)$ the set of finite positive Radon measures on $\pG$. Then $\Gamma$ acts on $\clM(\pG)$ via $g_*\nu(A) = \nu(g^{-1}A)$ for any Borel set $A$ in $\pG$.

\begin{defi}
We call a map $x \mapsto \nu_x$ equivariant if for every $x,g\in \Gamma$, we have
\[
\nu_{gx}(A) = \nu_{x}\left(g^{-1}A\right)
\]
for every Borel set $A \subset\pG$.
\end{defi}



\begin{defi}
Let $\omega \in [0, \infty\mathclose{[}$. We call a $\Gamma$-equivariant map
\[
\nu: \Gamma \to
\clM(\pG),\quad x \mapsto \nu_x
\]
an
\emph{$\omega$-dimensional quasi-conformal density} if for any $x, y \in \Gamma$ the following holds
\begin{equation}\label{cdensitydef}
\frac{d\nu_{x}}{d\nu_{y}}(\xi) \asymp \rme^{-\omega B_\xi (x, y)}K_{\xi}(x,y|r),
\end{equation}
for $\nu_y$-a.e. points $\xi \in \pG$, where the implicit constant does neither depend on $x, y$, nor on $\xi$.
\end{defi}



We prove the following result.
\begin{lemm}
Let $\{\nu_x\}_{x\in\Gamma}$ be a $\sigma$-dimensional quasi-conformal density on $\pG$. Then the support of any $\nu_x$ is $\pG$.
\end{lemm}

\begin{proof}
By definition, the support $\supp(\nu_x)$ is a maximal closed subset such that any point in $\pG \setminus \supp(\nu_x)$ has an open neighborhood which is $\nu_x$-null. It is well-known that $\pG$ is a minimal $\Gamma$-invariant closed set, see for instance~\cite[Section~2]{Bo99}. Thus, it suffices to prove that the support of $\nu_x$ is $\Gamma$-invariant. This follows from equivariance and quasi-conformality, since $\nu_x(A)=\nu_{gx}(gA)$ for any Borel subset $A\subset \pG$ and $\nu_x$ and $\nu_{gx}$ are absolutely continuous with respect to each other.
\end{proof}

As explained in the introduction, we associate the following Poincaré series to $\mu$ and to the word distance, by setting
\begin{align*}
\Theta_\Gamma({r, s}) := \sum_{x\in \Gamma} G(e, x |r) \rme^{-s|x|} = \sum_{n\ge 0} H_r(n) \rme^{-sn} =G(e,e|r)\cdot \sum_{x\in \Gamma} \rme^{-s|x|-|x|_r}
\end{align*}
where we recall that $H_r(n)=\sum_{x\in S_n} G(e,x|r)$ and that the
\emph{critical exponent} $\omega_\Gamma(r)$ is defined by
\[
\omega_\Gamma(r) = \limsup\limits_{n \to \infty} \frac{1}{n} \log H_r(n).
\]
The group $\Gamma$ is of \emph{divergent} (resp.
\emph{convergent}) type {for the Green function} if $\Theta_\Gamma({r,s})$ is divergent (resp. convergent) at $s=\omega_\Gamma(r)$.

Recall $\fkd_r(x,y):=\omega_\Gamma(r)|x^{-1}y| +|x^{-1}y|_r$.
\begin{lemm}\label{NewPoincareSeries}
The series $\clP_\Gamma$ defined as
\[
\forall s>0,\;\clP_{\Gamma}(s):=\sum\limits_{x \in \Gamma} \rme^{- s \fkd_r(x,y)}
\]
has critical exponent $1$, and the divergence of $\clP_{\Gamma}(s)$ at $s=1$ is equivalent to that of $\Theta_\Gamma({r,s})$ at $s=\omega_\Gamma(r)$.
\end{lemm}

\begin{proof}
By~\cite[Lemma~2.1]{GouezelLalley}, $G(e,x|R)$ goes to 0 as $|x|$ tends to infinity. Thus, $G(e,x|r)\leq G(e,x|R)\leq 1$ for large enough $x$. In particular, we see that for $s> 1$,
\[
\sum_{x\in \Gamma}G(e,x|r)^s\rme^{-s\omega_\Gamma(r) |x|}\leq C \sum_{x\in \Gamma}G(e,x|r)\rme^{-s\omega_\Gamma(r) |x|}.
\]
The right-hand side converges, so the left-hand side also converges. Conversely, for $s<1$,
\[
\sum_{x\in \Gamma}G(e,x|r)\rme^{-s\omega_\Gamma(r) |x|}\leq C \sum_{x\in \Gamma}G(e,x|r)^s\rme^{-s\omega_\Gamma(r) |x|}.
\]
The left-hand side diverges, so the right-hand side also diverges. Thus, the critical exponent of $\clP_\Gamma(s)$ is 1.

Also, note that $\clP_\Gamma(1)=\Theta_\Gamma(r,\omega_\Gamma(r))$, so the second conclusion of the lemma follows.
\end{proof}

Write $\mu(f) = \int f d\mu$ for a continuous function $f \in C(\pG)$. We endow $\clM(\pG)$ with the weak topology. That is, a sequence $\mu_n \in \clM(\pG)$ converges to $\mu$ if $\mu_n(f) $ converges to $ \mu(f)$ for any $f \in C(\pG)$. Equivalently, by the Portmanteau Theorem~\cite[Theorem~2.1]{Bill99}, $\mu_n$ converges to $\mu$ if $\liminf_{n \to \infty}\mu_n(U) \ge \mu(U)$ for any open set $U \subset \pG$.

We start by constructing a family of measures $\{\nu_x^s\}_{x\in \Gamma}$ supported on $\Gamma$ for any $s >1$. First, assume that $\clP_\Gamma(s)$ is divergent at $s=1$. Set
\[
\nu^s_x = \frac{1}{\clP_{\Gamma}(s)} \sum\limits_{z \in \Gamma} \rme^{-s\fkd_r(x,z)} \cdot \dirac{z},
\]
where $s >1$ and $x \in \Gamma$. Note that $\nu^s_x$ is a probability measure.


On the contrary, assume that $\clP_{\Gamma}(s)$ is convergent at $s=1$, Patterson introduced in~\cite[Lemma~3.1]{Patt} a monotonically increasing function $H: \R_{\ge 0} \to \R_{\ge 0}$ with the following property:
\begin{equation}\label{patterson}
\forall \epsilon >0, \;\exists t_\epsilon,\quad \forall t > t_\epsilon,\quad \forall a>0,\quad  H(a+t) \le \exp(a\epsilon) H(t).
\end{equation}
and such that the following modified series
\[
\clP_{\Gamma}'(s):=\sum\limits_{x \in \Gamma} H(\fkd_r(x,z)) \cdot \rme^{-s\fkd_r(x,z)}
\]
is divergent for $s \le 1$ and convergent for $s>1$.
Then define measures as follows:
\[
\nu^s_x = \frac{1}{\clP_{\Gamma}'(s)} \sum\limits_{z \in \Gamma} \rme^{-s \fkd_r(x,z)} \cdot H(\fkd_r(x,z)) \cdot \dirac{z},
\]
where $s >1$ and $x \in \Gamma$.

Choose $s_i \searrow 1$ such that $\nu_x^{s_i}$ are convergent in $\clM(\pG\cup\Gamma)$ for all $x\in \Gamma$. Let $\nu_x = \lim \nu_x^{s_i}$ be the limit measures, which are called \emph{Patterson--Sullivan measures} associated with the Poincaré series $\clP_\Gamma$. Note that forcing the Poincaré series to be divergent at 1, we have $\nu_x(\pG) = 1$. In the sequel, we write PS-measures as shorthand for Patterson--Sullivan measures. We also write $\partial\Gamma^{\con}$ for the set of conical limit points in the Bowditch boundary.

\begin{prop}\label{prequasiconf}
The PS-measures $\{\nu_x\}_{x\in \Gamma}$ on the Bowditch boundary are absolutely continuous with respect to each other and satisfy
\begin{align}
\forall \xi\in \pG ,\quad  \frac{d\nu_x}{d\nu_y}(\xi) &\ge \rme^{-\fkd_r(x,y)},\label{LowerBoundEQ}\\
\forall \nu_y\; \text{a.e.}\; \xi\in \pG^{\con},\quad  \frac{d\nu_{x}}{d\nu_{y}}(\xi) &\asymp \rme^{-\omega_\Gamma(r) B_\xi (x, y)}K_{\xi}(x,y|r),\label{cdensity}
\end{align}
where the implicit constant is independent of $x,y$ and $\xi$.
\end{prop}


\begin{rema}
If $\Gamma$ is of divergent type for Green function, then Theorem~\href{https://ahl.centre-mersenne.org/articles/10.5802/ahl.256/}{4.17} below says that PS-measures have no atoms on Bowditch boundary and give full measure to conical limit points, so~\eqref{cdensity} holds for $\nu_y$-a.e. $\xi\in \pG$. In this case, $\nu$ is an $\omega_\Gamma(r)$-dimensional quasi-conformal density.
\end{rema}

\begin{proof}
Since $\fkd_r$ satisfies the triangle inequality and $\lim_{t\to\infty} \frac{H(a+t)}{H(t)}=1$, we see that $\{\nu_x: x\in \Gamma\}$ are absolutely continuous with respect to each other,
\begin{align*}
\rme^{-\fkd_r(x,y)}\le \frac{d\nu_x}{d\nu_y}(\xi) \le \rme^{\fkd_r(x,y)}.
\end{align*}


We now verify the quasi-conformality at conical limit points. We only consider the case where $\Gamma$ is of convergent type, the divergent type being simpler. Let $\epsilon> 0$ and $t_\epsilon$ the number be given by~\eqref{patterson} for the function $H$. Let $\phi=\frac{d\nu_x}{d\nu_y}$ be the Radon--Nikodym derivative, uniquely defined up to a $\nu_y$-null set. Let $\xi \in \pG$ be a conical limit point and consider the open neighborhood $V$ of $\xi$ and the uniform constant $C$ given by Lemma~\ref{BusemanEstLem}.

Let $f$ be a continuous function supported in $V$. One can choose $V$ also such that $\fkd_r(y, z) > t_\epsilon$ for all $z \in V$. If $z \in V$ satisfies $\fkd_r(x, z) >\fkd_r(y, z)$, then we have
\begin{align*}
H(\fkd_r(x,z)) &= H(\fkd_r(x,z)-\fkd_r(y,z)+\fkd_r(y,z)) \\
& \le \rme^{\epsilon[\fkd_r(x,z)-\fkd_r(y,z)]} \cdot H(\fkd_r(y,z)) \\
& \le \rme^{\epsilon (C + \fkB_\xi(x, y))}\cdot H(\fkd_r(y, z)).
\end{align*}
Since $H$ is increasing, we have
\begin{equation}\label{prequasiconfeq}
C_\epsilon^{-1} H(\fkd_r(y, z))\le H(\fkd_r(x, z)) \le C_\epsilon H(\fkd_r(y, z)),
\end{equation}
where $C_\epsilon= \rme^{\epsilon (C + \fkB_\xi(x, y))}>1$ depends on $\epsilon, C$ and $ (x,y)$, but not on $z\in V$. Note that $C_\epsilon\to 1$ as $\epsilon\to 0$. By symmetry, the conclusion~\eqref{prequasiconfeq} also holds if $\fkd_r(x, z) < \fkd_r(y, z)$ for $z \in V$.

Using~\eqref{prequasiconfeq}, we get the following estimates. First,
\begin{align*}
\nu^s_x(f) & = \frac{1}{\clP_{\Gamma}(s)} \sum\limits_{z \in V}
\rme^{-s \fkd_r(x,z)} H(\fkd_r(x,z)) {f(z)} \\
& \leq C_\epsilon \rme^{-s\fkB_{z}(x, y))}\cdot \frac{1}{\clP_{\Gamma}(s)}
\sum\limits_{z \in V}
\rme^{-s\fkd_r(y, z)} H(\fkd_r(y, z)) {f(z)}\\
& \asymp C_\epsilon \rme^{-s\fkB_\xi(x, y)} \nu^s_y(f)
\end{align*}
and second
\begin{align*}
\nu^s_x(f) & = \frac{1}{\clP_{\Gamma}(s)} \sum\limits_{z \in V}
\rme^{-s \fkd_r(x,z)} H(\fkd_r(x,z)) {f(z)} \\
& \geq C_\epsilon^{-1} \rme^{-s\fkB_{z}(x, y))}\cdot \frac{1}{\clP_{\Gamma}(s)}
\sum\limits_{z \in V}
\rme^{-s\fkd_r(y, z)} H(\fkd_r(y, z)) {f(z)}\\
& \asymp C_\epsilon \rme^{-s\fkB_\xi(x, y)} \nu^s_y(f)
\end{align*}
where the implicit constants depend only on $C$ but not on $\epsilon$. Letting $s\to 1$, we get
\[
C_\epsilon^{-1} \rme^{-\fkB_\xi(x, y)}\nu_y(f)\prec \nu_x(f) \prec C_\epsilon \rme^{-\fkB_\xi(x, y)} \nu_y(f)
\]
for any continuous function $f$ supported in $V$. As $\epsilon\to 0$, $C_\epsilon\to 1$, hence it follows that
\[
\phi(\xi) \asymp \rme^{-
\fkB_\xi (x, y)}
\]
for $\nu_y$-a.e. conical limit point $\xi \in \pG$.
\end{proof}

\subsection{Shadow Lemma}

The key observation in the theory of PS-measures is the Sullivan Shadow lemma shedding some light on the relation between the PS-measure and the geometric properties of the boundary.

Recall that $\partial \Gamma$ is a visual boundary for the Cayley graph $\Cay(G,S)$: any two distinct points $x,y\in \Cay(G,S)\cup \partial \Gamma$ can be connected by a geodesic. See~\cite[Proposition~2.4]{GePoJEMS}.

\begin{defi}
Let $C>0$ and $x \in\Gamma$. The \emph{shadow $\Pi_C(x)$} at $x$ is the set of points $\xi \in \pG$ such that there exists some geodesic $\gamma=[e, \xi]$ intersecting $B(x, C)$.

The \emph{partial shadow $ \Psi_{C}(x)$} at $x$ is the set of points $\xi \in \pG$ such that some geodesic $[e,\xi]$ contains a transition point $C$-close to $x$.
\end{defi}

We now state the Shadow lemma in our context, whose proof follows closely the proofs of~\cite[Lemmas~4.1 \& 4.2]{YangPS} with Lemma~\ref{BusemanEstLem} replacing Lemma~2.19 there. Let us denote by $\Psi^{\con}_C(g)$ the set of all conical limit points in $\Psi_C(g)$.

\begin{lemm}[Shadow lemma]\label{ShadowLem}
Let $\{\nu_x\}_{x \in \Gamma}$ be an $\omega_\Gamma(r)$-dimensional PS measures on the Bowditch boundary $\pG$. Then there exists $C_0> 0$ such that for any $C\ge C_0$ and $x\in \Gamma$ the following two inequalities hold
\begin{align}
\rme^{-\omega_\Gamma(r) |x|}G(e,x|r) &\prec_C \, \nu_e(\Psi_{C}(x))\le \nu_e(\Pi_{C}(x)),\label{LowBoundEQ}\\
\nu_e(\Psi^{\con}_{C}(x)) &\prec_C \, \rme^{-\omega_\Gamma(r) |x|} G(e,x|r).\label{UpperBoundEQ}
\end{align}
\end{lemm}

\begin{rema}
If $\nu_e$ has no atoms at parabolic points which form a countable subset of the Bowditch boundary, then we obtain the full strength of the partial shadow lemma without having to restrict our attention to conical points. The upper bound~\eqref{UpperBoundEQ} for the partial shadow uses the relative Ancona inequalities (Lemma~\ref{weakAncona}), while it is unknown whether the upper bound holds for the usual shadow.
\end{rema}

\begin{proof}
Let $F$ be a set of three loxodromic elements with pairwise disjoint fixed points. For each $f\in F$, let $\alpha:=\bigcup_{i\in \Z} f^i[e,f]$ be an $\langle f\rangle$-invariant quasi-geodesic between two fixed points $f_-,f_+\in\pG$. Let $U_f\subset \pG$ be an open neighborhood of $f_+$ so that for any $\eta\in U_f$, the projection of $\eta$ to the axis $\alpha$ has a distance to $e$ at least $C$. By~\cite[Lemma~2.4]{DWY}, for any $x\in \Gamma$, there exists $f\in F$ so that $[x^{-1},\eta]$ contains a transition point $C$-close to $e$. Thus, $U_f\subset x^{-1}\Psi_{C}(x)$.

As $\Gamma$ acts minimally with a dense orbit in $\pG$, the $\Gamma$-orbit of any open set $U\subset \pG$ covers $\pG$, so $U$ have positive $\nu_e$-measure. Hence, setting
\[
D=\min\{\nu_e(U_f): f\in F\}>0
\]
which is independent of $x$, we have
\[
\nu_e\left(x^{-1}\Psi_{C}(x)\right) \ge D.
\]
Since $\nu_x $ is equivariant, the lower bound in~\eqref{LowerBoundEQ} implies
\begin{align*}
\nu_e(\Psi_{C}(x))={\nu_{x^{-1}}\left(x^{-1}\Psi_{C}(x)\right)}&\ge \rme^{-\omega_\Gamma(r) |x|}\frac{G(e,x|r) }{G(e,e|r)}\cdot {\nu_e\left(x^{-1}\Psi_C(x)\right)}\\
&\ge D \rme^{-\omega_\Gamma(r) |x|}\frac{G(e,x|r)}{G(e,e|r)}.
\end{align*}
Since $G(e,e|r)$ is bounded by $G(e,e|R)$, this concludes the proof of the lower bound.

For any $\xi \in x^{-1} \Psi^{\con}_C(x)$, there is a geodesic $\gamma$ from $x^{-1}$ to $\xi$ which intersects $B(e, C)$ and contains a transition point. Thus, $|B_\xi(x^{-1},e)
- d(x^{-1}, e)| \le 2C.$ By the relative Ancona inequalities (Lemma~\ref{weakAncona}), there is a constant $C_1$ independent of $r$ such that
\[
K_\xi\left(x^{-1}, e|r\right) = \lim_{z\to \xi}\frac{G\left(x^{-1},z|r\right)}{G(e,z|r)}
\le C_1 G(e,x|r).
\]
Also, by~\eqref{cdensity} there is a constant $C_2>0$ such that for $\nu_e$-a.e. conical limit point $\xi \in \pG^{\con}$,
\[
\frac{d\nu_{x^{-1}}}{d\nu_e}(\xi) \le C_2 \rme^{-\omega_\Gamma(r) B_\xi\left(x^{-1},e\right)} K_\xi\left(x^{-1}, e| r\right).
\]
Combining together the above estimates, we have
\begin{align*}
\nu_e(\Psi^{\con}_C(x)) & = \nu_{x^{-1}}\left(x^{-1}\Psi^{\con}_C(x)\right) \\
& \le C_2\int_{x^{-1}\Psi^{\con}_C(x)} \rme^{-\omega_\Gamma(r) B_\xi\left(x^{-1},e\right)} K_\xi\left(x^{-1}, e|r\right) d \nu_e(\xi) \\
& \leq C_1C_2 \rme^{2C\omega_\Gamma(r)} \cdot \rme^{-\omega_\Gamma(r) |x|} G(e,x|r),
\end{align*}
which finishes the proof of the upper bound.
\end{proof}

\begin{thebibliography}{99}
\bibitem[Anc88]{Ancona} Alano Ancona, ``Positive harmonic functions and hyperbolicity'', \emph{Potential theory-surveys and problems}, Lecture Notes in Mathematics1344, Springer,1988,1--23.
\bibitem[Bil99]{Bill99} Patrick P.Billingsley, \emph{Convergence of probability measure},2nd ed., John Wiley \& Sons,1999.
\bibitem[BB07]{BlachereBrofferio} Sébastien Blachère and Sara Brofferio, ``Internal Diffusion Limited Aggregation on discrete groups having exponential growth'', \emph{Probab. Theory Relat. Fields}\textbf{137}(2007),no.3-4,323--343.
\bibitem[Bow99]{Bo99} Brian H.Bowditch, ``Convergence groups and configuration spaces'', \emph{Geometric group theory down under (Canberra,1996)}, Walter de Gruyter,1999,23--54.
\bibitem[Bow12]{Bowditch} Brian H.Bowditch, ``Relatively hyperbolic group'', \emph{Int. J. Algebra Comput.}\textbf{22}(2012),no.3, article no.1250016 (66pages).
\bibitem[DG21]{DG21} Matthieu Dussaule and Ilya Gekhtman, ``Stability phenomena for Martin boundaries of relatively hyperbolic groups'', \emph{Probab. Theory Relat. Fields}\textbf{179}(2021),no.1-2,201--259.
\bibitem[DWY25]{DWY} Matthieu Dussaule, Longmin Wang and Wenyuan Yang, ``Branching random walks on relatively hyperbolic groups'', \emph{Ann. Probab.}\textbf{53}(2025),no.2,391--452.
\bibitem[GP15]{GePoGGD} Victor Gerasimov and Leonid Potyagailo, ``Non-finitely generated relatively hyperbolic groups and Floyd quasiconvexity'', \emph{Groups Geom. Dyn.}\textbf{9}(2015),no.2,369--434.
\bibitem[GP13]{GePoJEMS} Victor Gerasimov and Leonid Potyagailo, ``Quasi-isometric maps and Floyd boundaries of relatively hyperbolic groups'', \emph{J. Eur. Math. Soc.}\textbf{15}(2013),no.6,2115--2137.
\bibitem[GGPY21]{GGPY} Ilya Gekhtman, Victor Gerasimov, Leonid Potyagailo and Wenyuan Yang, ``Martin boundary covers Floyd boundary'', \emph{Invent. Math.}\textbf{223}(2021),no.2,759--809.
\bibitem[GP16]{GP16} Victor Gerasimov and Leonid Potyagailo, ``Quasiconvexity in relatively hyperbolic groups'', \emph{J. Reine Angew. Math.}\textbf{710}(2016),95--135.
\bibitem[Gou14]{Gouezel} Sébastien Gouëzel, ``Local limit theorem for symmetric random walks in Gromov-hyperbolic groups'', \emph{J. Am. Math. Soc.}\textbf{27}(2014),no.3,893--928.
\bibitem[GL13]{GouezelLalley} Sébastien Gouëzel and Steven P.Lalley, ``Random walks on co-compact Fuchsian groups'', \emph{Ann. Sci. Éc. Norm. Supér.}(4)\textbf{46}(2013),no.1,131--175.
\bibitem[Pat76]{Patt} Samuel J.Patterson, ``The limit set of a Fuchsian group'', \emph{Acta Math.}\textbf{136}(1976),no.3-4,241--273.
\bibitem[Yan22]{YangPS} Wenyuan Yang, ``Patterson--Sullivan measures and growth of relatively hyperbolic groups'', \emph{Peking Math. J.}\textbf{5}(2022),153--212.
\end{thebibliography}
\end{document}
